\documentclass[12pt]{ximera}
\title{Homework 5: Conditional Probability and Independence}
\begin{document}
\begin{abstract}
Section 7.5: conditional probability from the definition, testing independence with
real injury data, and a full tree diagram for two draws without replacement.
\end{abstract}
\maketitle

\section*{Section 7.5: Conditional Probability and Independence}

\begin{problem}
Use the formulas for conditional probability to solve the following. Give exact
answers.

Suppose $P(E)=0.6$, $P(F)=0.4$, and $P(E\cap F)=0.15$.

\begin{prompt}
\textbf{(a)} What is $P(F\mid E)$?

$P(F\mid E) = \answer[tolerance=0.001]{0.25}$

\begin{hint}
$P(F\mid E)=\dfrac{P(E\cap F)}{P(E)}$.
\end{hint}

\begin{feedback}
\[
P(F\mid E)=\frac{P(E\cap F)}{P(E)}=\frac{0.15}{0.6}=0.25.
\]
\end{feedback}
\end{prompt}

\begin{prompt}
\textbf{(b)} Are the events $E$ and $F$ independent?

\begin{multipleChoice}
\choice{Yes}
\choice[correct]{No}
\end{multipleChoice}

\begin{feedback}
No. Independence would require $P(F\mid E)=P(F)$, but $P(F\mid E)=0.25$ while
$P(F)=0.4$.

Equivalently, independence would require $P(E\cap F)=P(E)\cdot P(F)$, but
$0.6\times 0.4 = 0.24 \neq 0.15$. Knowing $E$ happened makes $F$ less likely, so the
two events are dependent.
\end{feedback}
\end{prompt}
\end{problem}

\begin{problem}
The following table lists injuries for ice hockey players wearing either a half shield
or a full shield, as reported in a university athletics study.

\begin{center}
\begin{tabular}{r|cc|c}
 & Half Shield ($H$) & Full Shield ($F$) & Totals\\\hline
Head and Face Injuries ($A$) & 95 & 34 & 129\\
Concussions ($B$) & 41 & 38 & 79\\
Neck Injuries ($C$) & 9 & 7 & 16\\
Other Injuries ($D$) & 202 & 150 & 352\\\hline
Totals & 347 & 229 & 576
\end{tabular}
\end{center}

\begin{prompt}
\textbf{(a)} What is the probability that a player who sustained a head or face injury
was wearing a half shield? Write it in conditional notation, then solve.

The probability to find is
\begin{multipleChoice}
\choice[correct]{$P(H\mid A)$}
\choice{$P(A\mid H)$}
\choice{$P(A\cap H)$}
\end{multipleChoice}

As a fraction: $\answer{95/129}$

\begin{feedback}
We are told the injury was a head or face injury, so we restrict to row $A$: this is
$P(H\mid A)$. Of the $129$ head and face injuries, $95$ involved a half shield:
\[ P(H\mid A)=\frac{95}{129}\approx 0.736. \]
\end{feedback}
\end{prompt}

\begin{prompt}
\textbf{(b)} Was the risk of head and face injury independent of wearing a half shield
in these cases?

\begin{multipleChoice}
\choice{Yes, independent}
\choice[correct]{No, not independent}
\end{multipleChoice}

\begin{feedback}
Compare $P(H\mid A)$ with $P(H)$. Overall, $P(H)=\frac{347}{576}\approx 0.602$, but among
head and face injuries $P(H\mid A)\approx 0.736$. Since these differ, the events are not
independent: half shields are over-represented among head and face injuries.
\end{feedback}
\end{prompt}
\end{problem}

\begin{problem}
Suppose you have a container of assorted chocolate truffles with $3$ milk chocolates
(M), only one dark chocolate (D), and $5$ white chocolates (W) remaining --- nine
truffles in all.

You draw one for yourself and then one to share with a friend, \emph{without}
replacing the first. Below, the truffle drawn for you is labeled with a $1$, and
your friend's with a $2$.

Enter every answer in this problem as a fraction.

\begin{prompt}
\textbf{(a, first draw)} The probabilities on the first three branches of the tree:

$P(M_1)=\answer{1/3}$ \qquad
$P(D_1)=\answer{1/9}$ \qquad
$P(W_1)=\answer{5/9}$

\begin{hint}
There are $9$ truffles at the first draw.
\end{hint}

\begin{feedback}
Out of $9$ truffles, $3$ are milk, $1$ is dark, and $5$ are white:
\[
P(M_1)=\frac39=\frac13, \qquad P(D_1)=\frac19, \qquad P(W_1)=\frac59.
\]
These add to $1$, as the first-generation branches of a tree always must.
\end{feedback}
\end{prompt}

\begin{prompt}
\textbf{(a, second draw)} Now the conditional probabilities on the second set of
branches. Remember that one truffle has already been removed, so only $8$ remain.

Given your truffle was milk: \\
$P(M_2\mid M_1)=\answer{1/4}$ \quad
$P(D_2\mid M_1)=\answer{1/8}$ \quad
$P(W_2\mid M_1)=\answer{5/8}$

\smallskip
Given your truffle was dark: \\
$P(M_2\mid D_1)=\answer{3/8}$ \quad
$P(D_2\mid D_1)=\answer{0}$ \quad
$P(W_2\mid D_1)=\answer{5/8}$

\smallskip
Given your truffle was white: \\
$P(M_2\mid W_1)=\answer{3/8}$ \quad
$P(D_2\mid W_1)=\answer{1/8}$ \quad
$P(W_2\mid W_1)=\answer{1/2}$

\begin{hint}
In each row, subtract one from the count of the flavor you already drew, and use $8$
as the denominator. There was only one dark truffle to begin with.
\end{hint}

\begin{feedback}
After your draw, $8$ truffles remain, with one fewer of whichever flavor you took.
\begin{itemize}
\item If you took milk, the remaining pool is $2$ M, $1$ D, $5$ W:
      $\frac28=\frac14$, $\frac18$, $\frac58$.
\item If you took the dark one, the pool is $3$ M, $0$ D, $5$ W:
      $\frac38$, $0$, $\frac58$. The dark branch is impossible because there was
      only ever one dark truffle.
\item If you took white, the pool is $3$ M, $1$ D, $4$ W:
      $\frac38$, $\frac18$, $\frac48=\frac12$.
\end{itemize}
\end{feedback}
\end{prompt}

\begin{prompt}
\textbf{(b)} Multiply along each path to get the probability of each outcome.

$P(M_1\cap M_2)=\answer{1/12}$ \quad
$P(M_1\cap D_2)=\answer{1/24}$ \quad
$P(M_1\cap W_2)=\answer{5/24}$

\smallskip
$P(D_1\cap M_2)=\answer{1/24}$ \quad
$P(D_1\cap D_2)=\answer{0}$ \quad
$P(D_1\cap W_2)=\answer{5/72}$

\smallskip
$P(W_1\cap M_2)=\answer{5/24}$ \quad
$P(W_1\cap D_2)=\answer{5/72}$ \quad
$P(W_1\cap W_2)=\answer{5/18}$

\begin{hint}
$P(X_1\cap Y_2)=P(X_1)\cdot P(Y_2\mid X_1)$ --- multiply the two numbers along the
path from the start of the tree to that outcome.
\end{hint}

\begin{feedback}
Multiplying along each path (all nine have denominator $72$ before reducing):
\[
\begin{array}{lll}
P(M_1\cap M_2)=\frac39\cdot\frac28=\frac{6}{72}=\frac{1}{12} &
P(M_1\cap D_2)=\frac39\cdot\frac18=\frac{3}{72}=\frac{1}{24} &
P(M_1\cap W_2)=\frac39\cdot\frac58=\frac{15}{72}=\frac{5}{24}\\[4pt]
P(D_1\cap M_2)=\frac19\cdot\frac38=\frac{3}{72}=\frac{1}{24} &
P(D_1\cap D_2)=\frac19\cdot 0=0 &
P(D_1\cap W_2)=\frac19\cdot\frac58=\frac{5}{72}\\[4pt]
P(W_1\cap M_2)=\frac59\cdot\frac38=\frac{15}{72}=\frac{5}{24} &
P(W_1\cap D_2)=\frac59\cdot\frac18=\frac{5}{72} &
P(W_1\cap W_2)=\frac59\cdot\frac48=\frac{20}{72}=\frac{5}{18}
\end{array}
\]
As a check, the nine numerators over $72$ add up to
$6+3+15+3+0+5+15+5+20=72$, so the outcomes total probability $1$.
\end{feedback}
\end{prompt}

\begin{prompt}
\textbf{(c)} What is the probability that you and your friend draw the same flavor of
truffle?

$\answer{13/36}$

\begin{hint}
Three of the nine outcomes have both truffles the same flavor. Add those three
probabilities.
\end{hint}

\begin{feedback}
The ``same flavor'' outcomes are $M_1\cap M_2$, $D_1\cap D_2$, and $W_1\cap W_2$:
\[
\frac{6}{72} + 0 + \frac{20}{72} = \frac{26}{72} = \frac{13}{36} \approx 0.361.
\]
Note that $D_1\cap D_2$ contributes nothing, since there is only one dark truffle.
\end{feedback}
\end{prompt}
\end{problem}

\end{document}
